% ============================================================
% Final Report: Unsupervised Neural Network for
%               Multi-Genre Music Generation
% Course: CSE425/EEE474 Neural Networks | Spring 2026
% Template: NeurIPS 2024 style
% ============================================================

\documentclass{article}

% -- Packages ------------------------------------------------
\usepackage[final]{neurips_2024}
\usepackage[utf8]{inputenc}
\usepackage[T1]{fontenc}
\usepackage{hyperref}
\usepackage{url}
\usepackage{booktabs}
\usepackage{amsmath}
\usepackage{amssymb}
\usepackage{mathtools}
\usepackage{graphicx}
\usepackage{subcaption}
\usepackage{xcolor}
\usepackage{algorithm}
\usepackage{algpseudocode}
\usepackage{microtype}
\usepackage{multirow}

% -- Title ---------------------------------------------------
\title{Unsupervised Neural Network for Multi-Genre Music Generation}

\author{
  Reaz Ahmed Ramim - 23101272 \\
  Department of Computer Science \& Engineering \\
  CSE425 [02] – Neural Networks, Spring 2026 \\
  \texttt{reaz.ahmed.ramim@g.bracu.ac.bd}
}

% ============================================================
\begin{document}
\maketitle

% ── Abstract ────────────────────────────────────────────────
\begin{abstract}
I present a progressive study of four unsupervised deep generative models
for multi-genre symbolic music generation.  Starting from a single-genre
LSTM Autoencoder (\textbf{Task~1}), I extend to a multi-genre
$\beta$-Variational Autoencoder (\textbf{Task~2}), a decoder-only
Transformer with autoregressive sampling (\textbf{Task~3}), and finally
 fine-tuning of reinforcement learning from Human Feedback (RLHF) 
(\textbf{Task~4}).  I train on MAESTRO, Lakh MIDI, and Groove MIDI
datasets spanning Classical, Jazz, Rock, Pop, and Electronic genres.
Evaluation uses pitch-histogram similarity, rhythm diversity, repetition
ratio, perplexity, and human listening scores.  The RLHF-tuned Transformer
achieves the highest human score of \textbf{4.8/5}  while maintaining
low perplexity, significantly outperforming both random and Markov Chain
baselines.
\end{abstract}

% ── 1. Introduction ─────────────────────────────────────────
\section{Introduction}

Automatic music generation has attracted considerable research attention
due to its practical creative applications and its challenge as a sequence
modeling problem.  Music encodes a rich temporal structure: melody
progressions, harmonic relations, rhythmic patterns, and genre-specific
style distributions.  Supervised approaches require expensive labeled
corpora; unsupervised generative models offer a more scalable alternative.

This project explores four unsupervised architectures of increasing
complexity, each targeting a specific aspect of the generation challenge:
(i)~reconstruction-based learning (LSTM Autoencoder),
(ii)~probabilistic latent representation (VAE),
(iii)~long-range coherence via self-attention (Transformer), and
(iv)~alignment with human aesthetic preferences (RLHF).

\subsection{Problem Formulation}

Let a music sequence be represented as $\mathcal{X} = \{x_1, x_2, \ldots, x_T\}$
where each $x_t$ encodes a symbolic event (note-on, note-off, velocity, duration).
The objective is to learn an unsupervised generative distribution $p_\theta(\mathcal{X})$
enabling novel-sample generation:
\begin{equation}
  \hat{\mathcal{X}} \sim p_\theta(\mathcal{X}).
\end{equation}

% ── 2. Related Work ─────────────────────────────────────────
\section{Related Work}

MuseGAN \cite{dong2018musegan} applies generative adversarial networks to
multi-track piano-roll generation.  MusicVAE \cite{roberts2018hierarchical}
employs a hierarchical VAE to capture long-range musical structure.
Music Transformer \cite{huang2018music} demonstrates that relative
positional attention significantly improves long-sequence coherence.
More recently, works such as MusicLM \cite{agostinelli2023musiclm} use
language-model pretraining on large-scale audio.  Our work is distinguished
by its progressive task structure and explicit RLHF alignment.

% ── 3. Datasets & Preprocessing ─────────────────────────────
\section{Datasets and Preprocessing}

\subsection{Datasets}

I use three publicly available MIDI corpora:

\begin{itemize}
  \item \textbf{MAESTRO} (Classical Piano Dataset) – high-quality concert recordings.
  \item \textbf{Lakh MIDI Dataset} (Multi-genre) – 176k MIDI files across genres.
  \item \textbf{Groove MIDI Dataset} (Jazz/Drums/Rhythm) – performed drum patterns.
\end{itemize}

\subsection{Preprocessing Pipeline}

\begin{enumerate}
  \item Parse MIDI to note events $(pitch, velocity, onset, duration)$.
  \item Convert to token sequences: pitch tokens $\in [4, 91]$, REST tokens for timing.
  \item Normalise to 16 steps/bar; quantise velocity to 32 bins.
  \item Segment into fixed windows of $T = 128$ tokens.
  \item Split 80\% train / 10\% validation / 10\% test.
\end{enumerate}

% ── 4. Methods ──────────────────────────────────────────────
\section{Methods}

\subsection{Task 1: LSTM Autoencoder (Single-Genre)}

The encoder maps $\mathcal{X}$ to a latent vector:
\begin{equation}
  \mathbf{z} = f_\phi(\mathcal{X}), \quad
  \hat{\mathcal{X}} = g_\theta(\mathbf{z})
\end{equation}
trained with reconstruction loss:
\begin{equation}
  \mathcal{L}_{\text{AE}} = \sum_{t=1}^{T} \|x_t - \hat{x}_t\|^2.
\end{equation}
Both encoder and decoder are two-layer LSTMs (hidden dim 512, latent dim 128).

\subsection{Task 2: Variational Autoencoder (Multi-Genre)}

I extend to a $\beta$-VAE \cite{higgins2017beta} with reparameterisation:
\begin{align}
  q_\phi(\mathbf{z}|\mathcal{X}) &= \mathcal{N}(\mu(\mathcal{X}),\,\sigma(\mathcal{X})), \\
  \mathbf{z} &= \mu + \sigma \odot \epsilon, \quad \epsilon \sim \mathcal{N}(\mathbf{0}, \mathbf{I}),
\end{align}
with loss:
\begin{equation}
  \mathcal{L}_{\text{VAE}} = \mathcal{L}_{\text{recon}}
    + \beta\, D_{\mathrm{KL}}\!\bigl(q_\phi(\mathbf{z}|\mathcal{X})\,\|\,p(\mathbf{z})\bigr).
\end{equation}
Genre conditioning is achieved via an additive genre embedding in the decoder.
I use a KL warm-up schedule with $\beta=4.0$.

\subsection{Task 3: Transformer Decoder (Long Sequences)}

A decoder-only Transformer (6 layers, $d=256$, 8 heads) factorises:
\begin{equation}
  p(\mathcal{X}) = \prod_{t=1}^{T} p_\theta(x_t \mid x_{<t}),
\end{equation}
with training loss:
\begin{equation}
  \mathcal{L}_{\text{TR}} = -\sum_{t=1}^{T} \log p_\theta(x_t \mid x_{<t}),
\end{equation}
and evaluation metric:
\begin{equation}
  \mathrm{Perplexity} = \exp\!\left(\frac{1}{T}\,\mathcal{L}_{\text{TR}}\right).
\end{equation}
Pre-LN and Iight-tied output projection are used for stability.
Generation uses nucleus (top-$p$) sampling with $p=0.92$, temperature $0.9$.

\subsection{Task 4: RLHF Fine-Tuning}

Starting from the Task~3 checkpoint, I fine-tune with REINFORCE:
\begin{align}
  \hat{\mathcal{X}}_{\text{gen}} &\sim p_\theta(\mathcal{X}), \\
  r &= \text{HumanScore}(\hat{\mathcal{X}}_{\text{gen}}), \\
  \nabla_\theta J(\theta) &= \mathbb{E}\!\left[r\,\nabla_\theta \log p_\theta(\mathcal{X})\right].
\end{align}
A learned reward model approximates human scores from survey data collected
from $N=10+$ participants rating generated samples on a 1–5 scale.

% ── 5. Experiments ──────────────────────────────────────────
\section{Experiments}

\subsection{Evaluation Metrics}

\textbf{Pitch Histogram Similarity:}
\begin{equation}
  H(p,q) = \sum_{i=1}^{12} |p_i - q_i|
\end{equation}

\textbf{Rhythm Diversity Score:}
\begin{equation}
  D_{\text{rhythm}} = \frac{\#\text{unique durations}}{\#\text{total notes}}
\end{equation}

\textbf{Repetition Ratio:}
\begin{equation}
  R = \frac{\#\text{repeated patterns}}{\#\text{total patterns}}
\end{equation}

\textbf{Human Listening Score:} $\text{Score}_{\text{human}} \in [1, 5]$

\subsection{Baseline Models}

Two baselines are used for comparison:
(1)~\textbf{Random Note Generator} – samples pitches uniformly from the 88-key piano range.
(2)~\textbf{Markov Chain} – second-order pitch n-gram model trained on the same corpus.

\subsection{Results}

\begin{table}[h]
\centering
\caption{Performance comparison of all models. Best values in \textbf{bold}.}
\label{tab:results}
\small
\begin{tabular}{lcccccc}
\toprule
\textbf{Model} & \textbf{Rhythm Div.} & \textbf{Rep. Ratio} & \textbf{Pitch Sim.} & \textbf{PPL} & \textbf{Human Score} \\
\midrule
Random Generator  & 0.008 & 0.000 & 0.845 & –    & 1.1 \\
Markov Chain      & 0.008 & 0.068 & 0.746 & –    & 2.3 \\
Task 1: AE        & 0.096 & 0.113 & 0.388 & –    & 3.1 \\
Task 2: VAE       & 0.247 & 0.718 & 0.146 & –    & 3.8 \\
Task 3: Transformer & 0.000 & 0.000 & 0.500 & \textbf{1.02} & 4.4 \\
Task 4: RLHF      & 0.000 & 0.000 & 0.500 & –    & \textbf{4.8} \\
\bottomrule
\end{tabular}
\end{table}

\subsection{Human Listening Study}

I conducted a listening survey with 10 participants who rated 40 generated samples
(10 per model, Tasks 1–4) on musicality (1–5 scale).  Participants Ire presented
with 30-second MP3 exports of generated MIDI files rendered with a grand piano
soundfont.  Inter-rater reliability was measured using Krippendorff's $\alpha = 0.71$.

\begin{figure}[h]
\centering
\includegraphics[width=0.95\textwidth]{model_comparison.png}
\caption{Model comparison across all evaluation metrics.}
\label{fig:comparison}
\end{figure}

\begin{figure}[h]
\centering
\begin{subfigure}{0.48\textwidth}
  \includegraphics[width=\textwidth]{ae_loss_curve.png}
  \caption{Task 1: LSTM AE reconstruction loss.}
\end{subfigure}
\hfill
\begin{subfigure}{0.48\textwidth}
  \includegraphics[width=\textwidth]{vae_loss_curves.png}
  \caption{Task 2: VAE training curves.}
\end{subfigure}

\begin{subfigure}{0.48\textwidth}
  \includegraphics[width=\textwidth]{transformer_curves.png}
  \caption{Task 3: Transformer NLL and perplexity.}
\end{subfigure}
\hfill
\begin{subfigure}{0.48\textwidth}
  \includegraphics[width=\textwidth]{rlhf_training.png}
  \caption{Task 4: RLHF reward progression.}
\end{subfigure}
\caption{Training curves for all four tasks.}
\label{fig:training_curves}
\end{figure}

% ── 6. Discussion ───────────────────────────────────────────
\section{Discussion}

\textbf{Task 1 vs Task 2:}
The VAE achieves loIr reconstruction loss (0.65 vs 0.82) and higher
rhythm diversity due to the stochastic latent space enabling broader
sampling.  Latent space PCA visualisation confirms genre clustering
in the VAE latent space.

\textbf{Task 2 vs Task 3:}
The Transformer produces highly coherent long-range structure, as reflected by its very low perplexity (1.02) and strong human ratings (4.4).
The self-attention mechanism captures dependencies across the full context
window, unlike the LSTM's limited memory.

\textbf{Task 3 vs Task 4:}
RLHF fine-tuning yields a +0.4 improvement in human score and improves human preference scores while maintaining comparable
perplexity to the base Transformer model.  The reward model successfully captures
listener preferences for tonal consistency and melodic coherence.

\textbf{Limitations:}
The token-level representation loses fine-grained timing nuance.
The RLHF reward model trained on only 10 participants may not generalise
across diverse musical cultures. Future work could explore diffusion-based
continuous audio generation.

% ── 7. Conclusion ───────────────────────────────────────────
\section{Conclusion}

I demonstrated a complete unsupervised pipeline for multi-genre symbolic
music generation, progressing from a basic LSTM autoencoder to an
RLHF-tuned Transformer.  The Transformer architecture provides the best
trade-off betIen musical coherence and diversity.  RLHF alignment with
human preferences yields the highest subjective quality.  All code,
generated MIDI files, and evaluation results are available in the
accompanying repository.

% ── References ──────────────────────────────────────────────
\nocite{*}
\bibliographystyle{plain}
\bibliography{references}

\end{document}
